\section{Whole differential ideals under the genuine original Taylor map}

Write $R$ for the original analytic germ ring, $S=\mathbb C[[x_1,\ldots,x_n]]$,
and $T:R\to S$ for the genuine original Taylor homomorphism. The analytic
operators $D_i$ differentiate actual holomorphic representatives; the formal
operators $\partial_i$ are the actual formal power-series derivations.

\begin{definition}[Differential-generated ideals in the actual rings]
\label{hir:actual-original-differential-ideal-definition}
\uses{hir:actual-whole-original-taylor-coordinate-derivative}
\lean{Hironaka.HironakaActualOriginalDifferentialIdeals.originalDifferentialIdeal,Hironaka.HironakaActualOriginalDifferentialIdeals.formalDifferentialIdeal}
\leanok
For whole ideals $I\subset R$ and $K\subset S$, set
\[
 \mathcal D_R(I)=I+\langle D_i g:g\in I,\ i=1,\ldots,n\rangle_R,
 \qquad
 \mathcal D_S(K)=K+\langle\partial_iF:F\in K,\ i=1,\ldots,n\rangle_S.
\]
The original ideals are retained. These are ideals generated by derivatives
of all actual members, rather than derivatives of a supplied finite list.
\end{definition}

\begin{theorem}[Entire differential-generated ideals commute with actual Taylor extension]
\label{hir:actual-original-whole-differential-ideal-extension}
\uses{hir:actual-original-differential-ideal-definition,hir:actual-whole-original-taylor-coordinate-derivative}
\lean{Hironaka.HironakaActualOriginalDifferentialIdeals.map_originalDifferentialIdeal}
\leanok
For every original analytic ideal $I$,
\[
 T(\mathcal D_R(I))S=\mathcal D_S(T(I)S).
\]
In particular the right hand side includes derivatives of arbitrary formal
members of the whole extended ideal, with arbitrary formal coefficients.
\end{theorem}
\begin{proof}
Let $J=T(\mathcal D_R(I))S$ and $K=T(I)S$. Since $I\subset\mathcal D_R(I)$,
we have $K\subset J$. The genuine entire differentiation identity
$T(D_i g)=\partial_iT(g)$ shows that every generator of the extension of
$\mathcal D_R(I)$ lies in $\mathcal D_S(K)$, proving one inclusion.

For the reverse inclusion, prove $\partial_iF\in J$ for every $F\in K$
by induction on the literal ideal span of the actual image $T(I)$.
For a generator $F=T(g)$, whole Taylor differentiation identifies its
derivative with $T(D_i g)$, which lies in $J$. The zero and addition cases
follow from the genuine derivation laws and ideal closure. For an arbitrary
formal coefficient $a\in S$, the Leibniz identity gives
\[
 \partial_i(aF)=a\partial_iF+F\partial_i a.
\]
The first term belongs to $J$ by the induction hypothesis. The second does
because $F\in K\subset J$, with no requirement that $a$ converge or lift to
an analytic germ. Thus every genuine derivative generator of $\mathcal D_S(K)$
lies in $J$, as does $K$ itself. Both whole ideal inclusions prove equality.
\end{proof}

\begin{definition}[Every finite differential closure]
\label{hir:actual-original-iterated-differential-ideal-definition}
\uses{hir:actual-original-differential-ideal-definition}
\lean{Hironaka.HironakaActualOriginalDifferentialIdeals.originalDifferentialClosure,Hironaka.HironakaActualOriginalDifferentialIdeals.formalDifferentialClosure}
\leanok
Define $\mathcal D_R^0(I)=I$ and
$\mathcal D_R^{q+1}(I)=\mathcal D_R(\mathcal D_R^q(I))$;
define $\mathcal D_S^q(K)$ by the same actual natural-number recursion.
\end{definition}

\begin{theorem}[All actual finite differential closures commute with Taylor extension]
\label{hir:actual-original-all-differential-closure-extension}
\uses{hir:actual-original-whole-differential-ideal-extension,hir:actual-original-iterated-differential-ideal-definition}
\lean{Hironaka.HironakaActualOriginalDifferentialIdeals.map_originalDifferentialClosure}
\leanok
For every original analytic ideal $I$ and every $q\in\mathbb N$,
\[
 T(\mathcal D_R^q(I))S=\mathcal D_S^q(T(I)S).
\]
\end{theorem}
\begin{proof}
Induct on the actual integer $q$. At zero both sides are the same extension
of $I$. At the successor, apply the whole one-step theorem to the entire
original ideal $\mathcal D_R^q(I)$. Its right hand side is
$\mathcal D_S(T(\mathcal D_R^q(I))S)$, which the induction hypothesis
identifies with $\mathcal D_S(\mathcal D_S^q(T(I)S))$, exactly the recursive
successor. No finite truncation, output equality certificate or convergent
formal coefficient is assumed.
\end{proof}

These theorems concern the stated actual stalk ideals. Identifying them
with every source-facing derivative/Jacobian presentation, proving the
maximal Hilbert--Samuel cosupport criterion for general permissible centres,
and establishing transformed stability and Bennett inequalities remain
separate open tasks.
